\section{Limitations}
\label{sec:limitations}

\begin{itemize}
\item \textbf{Two data-center GPUs.} We measured speed on one NVIDIA L40S and
  one A100 (\S\ref{sec:a100}), never on a consumer GPU, the hardware this work
  targets. The persistent kernel ran on the L40S only.
\item \textbf{One calibration draw.} Every size is encoded from one draw of
  calibration text (\S\ref{sec:setup}). At 4B, with the earlier codebook, the
  spread across three draws exceeds every gain of Table~\ref{tab:chain}
  (\S\ref{sec:kept}). Those gains compare fixed files, but our gaps to FP16,
  AWQ and IQ2\_S do not include that spread.
\item \textbf{The scores and the speeds come from two forms of the engine.}
  MMLU and perplexity are scored on the dense reconstruction, and GSM8K on the
  per-matrix path, all with the expanded attention. The served engine is
  tied to them by a perplexity tolerance, by token identity and by a
  50-problem pilot per size (\S\ref{sec:setup}), not by a full run of each
  benchmark.
\item \textbf{Selection on the test set.} Our int4 choices used MMLU test
  questions (\S\ref{sec:chain}).
\item \textbf{GSM8K is an easy test for these models.} FP16 scores 92 to 95,
  and the problems have been public since 2021. We did not test Qwen3's
  thinking mode, which writes much longer chains, nor a harder math
  benchmark.
\item \textbf{Greedy decoding at batch 1.} The persistent kernel serves one
  request at a time and does not sample. Several requests go through a
  batched kernel that AWQ in vLLM outruns at every batch size
  (\S\ref{sec:context}).
\item \textbf{Long prompts.} After 8{,}192 prompt tokens our files hold 2.4 to
  3.6~GB more card memory than IQ2\_S in llama.cpp, and the 8B and 14B wait
  46 and 81~s for their first token (\S\ref{sec:context}).
\item \textbf{One llama.cpp build, one importance matrix.} The IQ2\_S files
  are ours, built from the calibration text of our files. Another build or
  another matrix may score differently.
\end{itemize}
